\subsection*{Annotation} A single annotator applied the frozen guideline, which
leaves inter-annotator agreement undefined. A blind second pass over
already-labelled items, shown in a different order with the first-pass labels
withheld, reproduced them exactly (27 of 27, $\kappa=1.00$). Self-agreement at
that level is consistent both with a guideline that admits one reading and with
an annotator recognising passages seen earlier, so the quantity we rest on is the
within-annotator context ablation (\S\ref{sec:models}), which measures how far the
same annotator's decisions depend on material outside the anchor sentence.
Whether two trained lawyers would agree on these labels is a question for a
follow-up study with a practitioner panel.

The annotator is a language model, which places it in the same family as some of
the systems under evaluation. Two design choices contain the circularity. The
supervised baselines are trained from the labels but are architecturally
independent of the label source, and they carry the context-scaling result. The
zero-shot baseline is a small open model external to the annotation process and
scored by label-token likelihood, so it is independent of the labels. We omit a
frontier model, which would report agreement with itself. The escalation results
in \S\ref{sec:long} rest on the same labels.

\subsection*{Benchmark and split sizes} Annotation produced \NBench{} usable labels,
\NDev{} for development and 104, 36, and 46 for the three held-out conditions. At
that size, differences of a few macro-$F_1$ points sit inside the noise, and the
temporal condition carries wide intervals. The two comparisons we rest on are the
per-stratum contrast, where the gap between the sparse model and the dictionary
on H1 is large, and the paired context comparison, which is within-item and
therefore better powered than the split sizes suggest. The separation between
supervised systems and the dictionary emerged as the development split grew,
which is the signature of an effect that needed data to surface and a reason to
rerun the benchmark at several thousand items. The issue-only control of
\S\ref{sec:models} indicates that part of the signal is subject matter, so later
versions should match strata on subject matter or report that control alongside
every headline number, as we do here.

\subsection*{Dictionary scope} A court can leave an issue open in ordinary prose.
Freezing the thirteen expressions before annotation protects against tuning them
to a result, and the recall audit in Appendix \ref{app:recall} bounds what they
miss by reading a random sample of opinions containing none of them. Prevalence
figures should therefore be read as the rate of explicitly marked non-resolution.
The item frame is scoped the same way, admitting anchors with a recoverable issue
statement, and Appendix \ref{app:issue} reports yield by stratum so the direction
of the loss stays visible.

\subsection*{Corpus coverage} CourtListener is broad but uneven. Coverage,
metadata completeness, and text quality vary by court and era, and the same
opinion is sometimes stored more than once from different upstream sources. We
deduplicate and report the eligibility funnel, and a change in measured
prevalence across periods remains partly a change in what was digitized and how,
which is why period-stratified rates appear throughout.

A citing passage is treated as being about the origin's issue when its content
words overlap the issue statement. That admits false positives, where a later
opinion cites the origin for a neighbouring point sharing vocabulary, and false
negatives, where a later court restates the issue in different words. Annotation
removes the false positives, since off-issue passages are labelled and dropped,
and the false-negative rate remains open.

\subsection*{Scope of the comparison} The comparison between unresolved and
decided issues matches on court, period, opinion length, and citation volume, and
is descriptive. Whether an issue was left open is a judicial choice correlated
with unobservables, including how hard and how contested the issue was, and those
plausibly affect how later courts describe it. We report the contrast as an
association.

Chains require the origin to carry a reporter citation locatable in the citing
text, and are restricted to citing opinions inside the population, so district
and state court treatments fall outside the sample and origins in the most recent
years have had little time to accumulate citations. Both restrictions bias the
chain sample toward older, better-reported, more-cited opinions.

The encoder baselines truncate at 512 tokens and the zero-shot model at 1600,
which caps the widest context condition, so the comparison across context widths
is partly a comparison of how much of the window each architecture uses. The
zero-shot model has 3.8B parameters, and its role here is to show that an
independent model responds to the context manipulation in the same direction. The
study covers English-language United States federal appellate opinions, and
extension to state courts, trial courts, agency adjudication, or other legal
systems remains future work.
